\documentclass[11pt]{article}
\usepackage[paperwidth=210mm,paperheight=5000mm,margin=20mm]{geometry}
\usepackage{amsmath,amssymb,xcolor,graphicx}
\usepackage[hypcap=false]{caption}
\usepackage[hidelinks,bookmarksnumbered,bookmarksopen]{hyperref}
\hypersetup{pdftitle={Linear algebra 7 — Eigenvalues and diagonalisation},pdfauthor={Lecturer: Dr A. Martin},
  pdfcreator={KherveNote}}
\renewcommand\labelitemii{\textbullet}
\renewcommand\labelitemiii{\textbullet}
\renewcommand\labelitemiv{\textbullet}
\renewcommand\labelenumii{\arabic{enumi}.\arabic{enumii}.}
\renewcommand\labelenumiii{\arabic{enumi}.\arabic{enumii}.\arabic{enumiii}.}
\renewcommand\labelenumiv{\arabic{enumi}.\arabic{enumii}.\arabic{enumiii}.\arabic{enumiv}.}
\setlength{\parindent}{0pt}
\setlength{\parskip}{0.6em}
\definecolor{knoteaccent}{HTML}{1A6DD8}
\definecolor{knotemuted}{HTML}{6B6B6B}
\definecolor{knotekey}{HTML}{D96B00}
\newcommand\knotetime[1]{\leavevmode\llap{\color{knotemuted}\scriptsize #1\hspace{1em}}}
\newcommand\knotekey[2][]{\par\noindent#1\colorbox{knotekey!12}{\parbox{\dimexpr\linewidth-2\fboxsep}{%
  \textbf{\color{knotekey}$\star$ Key point.} #2}}\par}
\newcommand\knotequestion[2][]{\par\noindent#1{\color{knoteaccent}\textbf{?}}~\emph{#2}\par}
\newenvironment{knotetranscript}{\par\color{knotemuted}\small}{\par}
\newcommand\knotesummary[1]{\par\noindent\fcolorbox{knoteaccent}{knoteaccent!6}{%
  \parbox{\dimexpr\linewidth-2\fboxsep-2\fboxrule}{\textbf{Summary}\par #1}}\par}
\makeatletter
\newdimen\knote@limit \knote@limit=4950mm
\newbox\knote@page \newbox\knote@chunk
\def\knote@ht#1{\dimexpr\ht#1+\dp#1\relax}
\def\knote@setheight#1{\global\pdfpageheight=#1\relax
  \special{pdf:pagesize width \the\paperwidth\space height \the\pdfpageheight}}
\def\knote@flush{\ifvoid\knote@page\else
  \knote@setheight{\dimexpr\knote@ht\knote@page+40mm+2mm\relax}%
  \box\knote@page\par\clearpage\global\pdfpageheight=\paperheight\fi}
\newenvironment{knotechunk}
  {\par\setbox\knote@chunk=\vbox\bgroup\hsize=\textwidth\linewidth=\textwidth}
  {\par\egroup
   \ifdim\knote@ht\knote@chunk>\knote@limit
     \knote@flush\unvbox\knote@chunk\par\clearpage
   \else\ifvoid\knote@page
     \global\setbox\knote@page=\vbox{\unvbox\knote@chunk}%
   \else\ifdim\dimexpr\knote@ht\knote@page+\knote@ht\knote@chunk\relax>\knote@limit
     \knote@flush\global\setbox\knote@page=\vbox{\unvbox\knote@chunk}%
   \else
     \global\setbox\knote@page=\vbox{\unvbox\knote@page\unvbox\knote@chunk}%
   \fi\fi\fi}
\AtEndDocument{\knote@flush}
\makeatother
\pagestyle{empty}
\begin{document}
\begin{knotechunk}
{\color{knoteaccent}\rule{\linewidth}{1.5pt}}\par
{\LARGE\bfseries Linear algebra 7 — Eigenvalues and diagonalisation}\par
{\large Lecturer: Dr A. Martin}\hfill{\color{knotemuted}05 October 2026 \textperiodcentered{} Maths building, Room 140}\par
{\color{knoteaccent}\rule{\linewidth}{0.6pt}}\par
\knotesummary{Eigenvectors are the directions a matrix only stretches; the eigenvalues are the stretch factors, found from det(A \ensuremath{-} \ensuremath{\lambda}I) = 0. A matrix with n independent eigenvectors can be written A = PDP\textsuperscript{-1}, which makes powers and many applications easy. Symmetric matrices always can, with an orthogonal P.}

\section{Eigenvalues and eigenvectors}

A non-zero vector $\mathbf{v}$ is an eigenvector of a square matrix $A$ when \[A\mathbf{v} = \lambda\mathbf{v}\] for some scalar $\lambda$, its eigenvalue: $A$ only stretches $\mathbf{v}$, it does not turn it.

\begin{itemize}\setlength{\itemsep}{0pt}\setlength{\parskip}{0pt}
  \item eigenvectors are only defined up to a scale factor
  \item $\lambda = 0$ is allowed; then $A$ is singular
  \item the eigenvalues of a triangular matrix are its diagonal entries
\end{itemize}

\knotekey{Sum of the eigenvalues = trace of $A$; product of the eigenvalues = $\det A$.}
\end{knotechunk}

\begin{knotechunk}
\section{The characteristic polynomial}

$(A - \lambda I)\mathbf{v} = 0$ has a non-zero solution only if \[\det(A - \lambda I) = 0\] — a polynomial of degree $n$ in $\lambda$, so an $n \times n$ matrix has $n$ eigenvalues counted with multiplicity (possibly complex).

\subsection{Worked example}

For $A = \begin{pmatrix} 2 & 1 \\ 1 & 2 \end{pmatrix}$: $\det(A-\lambda I) = (2-\lambda)^2 - 1 = (\lambda-3)(\lambda-1)$.

\begin{enumerate}\setlength{\itemsep}{0pt}\setlength{\parskip}{0pt}
  \item $\lambda_1 = 3$: $\mathbf{v}_1 = (1, 1)^T$
  \item $\lambda_2 = 1$: $\mathbf{v}_2 = (1, -1)^T$
\end{enumerate}

Check: trace $= 4 = 3 + 1$, det $= 3 = 3 \times 1$.

\knotequestion{What happens when an eigenvalue is repeated — are there always enough eigenvectors?}
\end{knotechunk}

\begin{knotechunk}
\section{Diagonalisation}

If $A$ has $n$ linearly independent eigenvectors, put them as the columns of $P$ and the eigenvalues on the diagonal of $D$: \[A = P D P^{-1}, \qquad A^k = P D^k P^{-1}\]

\begin{itemize}\setlength{\itemsep}{0pt}\setlength{\parskip}{0pt}
  \item distinct eigenvalues \ensuremath{\Rightarrow} independent eigenvectors \ensuremath{\Rightarrow} diagonalisable
  \item a repeated eigenvalue may give too few eigenvectors (a defective matrix, e.g. $\begin{pmatrix} 1 & 1 \\ 0 & 1 \end{pmatrix}$)
\end{itemize}

\knotekey{Spectral theorem: a real symmetric matrix has real eigenvalues and orthonormal eigenvectors, so $A = Q D Q^T$ with $Q^{-1} = Q^T$.}
\end{knotechunk}

\begin{knotechunk}
\section{Applications}

\begin{itemize}\setlength{\itemsep}{0pt}\setlength{\parskip}{0pt}
  \item Markov chains: the steady state is the eigenvector with $\lambda = 1$
  \item coupled oscillators: eigenvectors are the normal modes, eigenvalues give $\omega^2$
  \item principal component analysis: eigenvectors of the covariance matrix
  \item systems of ODEs $\dot{\mathbf{x}} = A\mathbf{x}$: solutions $e^{\lambda t}\mathbf{v}$
\end{itemize}

\knotequestion{Problem sheet 7, questions 3–6 for Friday.}
\end{knotechunk}
\end{document}
